% RECOVERY32: reconstructed from the immutable, fully reviewed cycle31 proof.
% These are new reconstruction bytes, not a claim to reproduce the lost32 hash.
\statement{A.95}{Proved}{A specified center-free sphere head and its invariant-free lifts}
Use the original marking
\((a,b,c,d,e,f)=(A_{12},A_{13},A_{23},A_{14},A_{24},A_{34})\).
Let H be a finite center-free group generated by specified elements
\(D_0,F_0\), and let V be a finite abelian H-module with \(V^H=0\).
Every homomorphism \(P_4\to V\rtimes H\) above the specified head
\[
 a\mapsto F_0,\quad c,d\mapsto D_0,\quad f\mapsto F_0,
 \quad b\mapsto(D_0F_0)^{-1},\quad e\mapsto(F_0D_0)^{-1}
 \tag{A.95.1}
\]
factors through that same marked free-two quotient. Consequently its
entire marked pentagon evaluation is an identity for every
\(u\in\widehat F(X,Y)'\) satisfying \(u(Y,X)=u(X,Y)^{-1}\).
No assumption on the characteristic or semisimplicity of V is needed.

\begin{proof}
Retain the actual inverse Artin action and full-twist identity of A.94:
\[
\begin{gathered}
 Z=abc,\quad P_3=F(a,c)\times\langle Z\rangle,
 \quad \operatorname{Ad}(Z^{-1})|_{F(d,e,f)}=\operatorname{Ad}(def),\\
 a^{-1}da=\operatorname{Ad}(de)d,\quad
 a^{-1}ea=\operatorname{Ad}(d)e,\quad [a,f]=1,\\
 [c,d]=1,\quad c^{-1}ec=\operatorname{Ad}(ef)e,
 \quad c^{-1}fc=\operatorname{Ad}(e)f.
\end{gathered}
\tag{A.95.2}
\]
These are the original action conventions already checked against
\citet[Example~1.1, pp.~465--466]{CECohen2010}; the pentagon word is
the marked display in \citet[p.~547, Eq.~(4)]{Furusho2010}.
Write \(A_0,B_0,C_0,D,E_0,F\) for the six images and set
\(P=DE_0F\), \(K=\langle D,E_0,F\rangle\). The image of K in H
is all H. A centralizer element of K has H-image in \(Z(H)=1\),
so belongs to V; there it centralizes K exactly when it belongs to
\(V^H\). Thus \(C_{V\rtimes H}(K)=1\).

The full-twist identity makes the images of Z and \(P^{-1}\) induce
the same conjugation on K, hence they are equal. The head of P is
\(D_0(F_0D_0)^{-1}F_0=1\), so P lies in V. Since Z is central
in the actual \(P_3\), the images of a and c commute with P.
Their heads generate H, hence \(P\in V^H=0\). Therefore P and Z
both have trivial image.

Now \(E_0=(FD)^{-1}\). Substitution into the complete a-row of
(A.95.2) makes \(\operatorname{Ad}(A_0^{-1})\) on K equal to
\(\operatorname{Ad}(F^{-1})\). Its trivial centralizer gives
\(A_0=F\). The complete c-row similarly gives \(C_0=D\).
Finally \(B_0=A_0^{-1}ZC_0^{-1}=(DF)^{-1}\). These are literal
equalities, not equalities only modulo a centralizer.

In the original pentagon
\[
 u(a,ce)u(bc,f)=u(c,f)u(ab,ef)u(a,c),
 \tag{A.95.3}
\]
the two left pairs become \((F,F^{-1})\), \((F^{-1},F)\),
and the middle right pair becomes \((D^{-1},D^{-1})\).
Their commutator-word values are one. The remaining equation is
\(1=u(D,F)u(F,D)\), which is precisely duality.

For a uniform example, let H be the dihedral group of order \(2n\)
with n odd and greater than one, generated by the prescribed
reflections. Over any characteristic-two finite field containing a
primitive nth root \(\zeta\), let its rotation act as
\(\operatorname{diag}(\zeta,\zeta^{-1})\) and a reflection swap
coordinates. The rotation has no invariant vector, and H is
center-free, so the theorem applies for every odd n.

The theorem concerns the specified head (A.95.1). It neither classifies
all nonabelian heads nor asserts that arbitrary module data are
canonical. A lift retaining a nontrivial boundary or a fiber
centralizer falls outside its hypotheses. No full cusp equation or
canonical-conjugacy conclusion follows.
\end{proof}
